\chapter{Conclusion}
\label{ch:conclusion}

Partially observed simultaneous-action games require agents to infer hidden
opponent state from interaction history and act without lookahead search. This
thesis studied whether a causal transformer with structured thirteen-token turns,
auxiliary belief supervision, and combinatorial held-out evaluation improves
offline imitation learning in that setting. Competitive Pok\'emon Gen~9 OU human
replays via Metamon served as the benchmark. The scientific contribution is
belief inference and combinatorial generalization under offline learning, not
competitive ladder play.

\section{Contributions}

The work delivers four concrete contributions.

\textbf{Pipeline and representation.} A reproducible path from Metamon replays to
encoded turn sequences, belief-enriched labels, and \texttt{poke-env} evaluation
hooks. Each turn is tokenized as thirteen discrete slots with causal history over
$N=8$ prior turns.

\textbf{Model family.} A shared \texttt{TurnTransformer} trunk with incremental
ablations: turn-only BC, matchup-conditioned BC
(\texttt{TeamMatchupEncoder}), and \texttt{ActiveBeliefTransformer} with move,
item, ability, and tera heads.

\textbf{Evaluation protocol.} Complementary splits for random battle holdout
(Split~A) and held-out opponent team keys (Split~B), with generalization gap as
the primary combinatorial statistic. A one-command scale script reproduces export,
training, and OOD evaluation.

\textbf{Empirical findings on human data.} Phase~1 BC reached 18.4\% action match
on $\approx$255\,000 turns (988 action classes). Phase~2 on 2\,000 battles showed
20.4\% random-holdout action match for turn-only BC, 28.1\% move-belief accuracy
for the belief-augmented model, 35.5\% for a latest-turn MLP baseline, and
near-zero OOD team gap after scaling. Explicit belief and matchup modules did not
improve policy metrics over turn-only cloning.

\section{Summary of Findings}

Belief inference and policy imitation responded differently to scale.
Move-belief accuracy more than doubled when the export grew from 500 to 2\,000
battles. Action accuracy improved for all models but did not rank belief-
augmented training above turn-only BC. Matchup conditioning over full team rosters
did not help, likely because species composition is already encoded in opponent
party tokens and because Gen~9 OU reveals all species at team preview.

The opponent-team generalization gap shrank from about two percentage points on
the pilot corpus to under half a point at scale for every architecture. Data
coverage therefore explains much of the early gap. Explicit belief and matchup
modules did not outperform the trunk on Split~B OOD accuracy.

Species belief, as originally envisioned for partially hidden rosters, does not
apply to this format. The empirical belief results concern active-opponent move
sets and related hidden attributes.

Offline imitation success should not be confused with live competitiveness.
Matching human actions on roughly one turn in five leaves most decisions wrong.
Against skilled human players the agent would be expected to lose almost every
game. That expectation does not erase the offline scientific findings above.

\section{Hypotheses Revisited}

\textbf{H1} received mixed support. The transformer learns history-conditioned
move beliefs well above pilot accuracy and beats a GRU baseline, but a
latest-turn MLP reaches higher move-belief accuracy (35.5\% versus 28.1\%).

\textbf{H2} was not supported at current scale. Auxiliary belief loss did not
improve action match over BC-only training, and belief fusion into the policy
head also failed to help.

\textbf{H3} was mixed. Scaling reduced OOD degradation for all models. Belief
and matchup modules did not achieve a smaller gap than the baseline.

\textbf{H4} (offline RL fine-tuning) remains future work.

The minimum viable bar in \Cref{ch:experiments} asked for at least two of four
criteria. This thesis meets the bar on (i)~demonstrating nontrivial belief
learning on hidden moves relative to chance and neural baselines and
(iv)~evidence that structured turn history and data scale contribute measurable
signal, but not on belief improving OOD policy or offline RL gains. The negative
result on belief-to-policy transfer is reported as a finding.

\section{Broader Impact}

The scientific contribution is not competitive Pok\'emon play. It is an empirical
study of transformer belief inference and combinatorial generalization under
offline imitation. The negative result on belief-to-policy transfer is informative
for domains where auxiliary prediction tasks are tempting but trunk capacity may
already absorb the relevant signal for action matching, or where multi-task
training fails to route belief features into the policy head. Practitioners
should validate whether explicit belief heads help policy metrics on their splits
before adding multi-task complexity.

\section{Future Work}

Priority directions follow from \Cref{ch:limitations,ch:analysis}.
\begin{itemize}
  \item Live win-rate evaluation on \texttt{poke-env} as a play-strength appendix,
    not as a replacement for the offline claims.
  \item Held-out set-combination splits (Split~C) and per-attribute belief tables
    for item, ability, and tera.
  \item Sweeps over $\lambda_\mathrm{bel}$ and larger models to re-test
    belief-to-policy transfer.
  \item Conservative offline RL (CQL or AMAGO) initialized from scaled BC
    checkpoints.
  \item Context-window and flat-sequence ablations, plus attention or
    decision-tree case studies on move-belief errors.
  \item Larger corpus scale toward the full Metamon subset used in Phase~1.
  \item Release of combinatorial split manifests as a standalone benchmark artifact.
\end{itemize}

Artifacts are available under \texttt{checkpoints/}, \texttt{logs/}, and
\texttt{scripts/scale\_belief\_research.py} for replication and extension.
